\scanpage{292}
\figure{\begin{tikzpicture}[x=1mm,y=1mm]
\draw[line cap=round,line width=.34mm] (0,3.4)--(.13,2.0);
\draw[line cap=round,line width=.37mm] (1.65,4.4)--(2.04,2.0);
\draw[line width=.40mm] (.48,.42) .. controls (.74,-.51) and (1.88,-.56) .. (2.11,.46);
\end{tikzpicture}}

On a un isom canonique

\[
T^p(M)=\operatorname{Hom}(T_p(M),\mathbb{Z}_p)
\]

On notera que a) $T^p$ n'est pas un foncteur exact, mais seulement exact ``mod groupes finis'' b) $T^p$ s'annule sur tous les groupes finis exactement c) De façon précise, $T^p$ établit une \inserted{\deleted{\unsure{anti}}}\unsure{équiv.|équi.} des catégories ``\uline{mod} les groupes finis'', i.e. essentiellement de \unsure{$\mathcal{C}$ mod finis|C-art finis} avec les espaces vectoriels de dim finie sur $\mathbb{Q}_p$.

Soit $M\unsure{\mathbin{-}|\to}N$ un hom de groupes de la forme $\mathbb{Q}_p/\mathbb{Z}_p^{h}$, et $H$ un $p$-groupe fini. \uline{Conditions équivalentes}

\[
T_p(N)\to T_p(M)\quad\text{injectif à conoyau libre [i.e. isom sur un facteur)}
\tag*{(i)}
\]

\[
T_p(N)\otimes H\to T_p(M)\otimes H\qquad\text{injectif}
\tag*{(i bis)}
\]

\[
\deleted{\operatorname{Tor}_1(\unsure{M|H},}\quad T_p(N)/pT_p(N)\to T_p(M)/pT_p(M)\qquad\text{injectif}
\tag*{(i ter)}
\]

\[
T^p(M)\to T^p(N)\qquad\text{surjectif}
\tag*{(ii)}
\]

\[
M\to N\qquad\text{épim sur un quotient facteur direct}
\tag*{(iii)}
\]

\[
{}_pM\to {}_pN\qquad\text{est \uline{surjectif}}
\tag*{(iv)}
\]

% archive sheet number: 291

